% !PRINTED 147
\weberpage{147}
\begin{center}
{\large \S.\ 45.}

\vspace{0.5em}
{\large\gesperrt{Second Proof of the Theorem on Symmetric Functions.}}
\end{center}

Let $S$ be a whole symmetric function of the variables $\alpha_1,\alpha_2\ldots\alpha_n$. The individual terms of this function are all of the form
\[
 M\alpha_1^{\nu_1}\alpha_2^{\nu_2}\ldots\alpha_n^{\nu_n},
\]
where $M$ is a coefficient independent of the $\alpha$. These terms are now to be arranged in a definite manner. After the order of the variables $\alpha_1,\alpha_2\ldots\alpha_n$ has been fixed, of two terms
\[
 A=M\alpha_1^{\nu_1}\alpha_2^{\nu_2}\ldots\alpha_n^{\nu_n},\qquad
 A'=M'\alpha_1^{\nu_1}\alpha_2^{\nu'_2}\ldots\alpha_n^{\nu'_n}
\]
% [sic] Source p147 omits the prime in the first exponent of A'; retained, not silently normalized.
$A$ is to be called the \gesperrt{higher} when the first of the differences\begingroup\renewcommand{\thefootnote}{E1}\footnote[1]{[Source note:] The first exponent of the second term is printed as $\nu_1$, without a prime. The following definition requires $\nu'_1$ there; the printed expression is retained unchanged.}\endgroup
\[
 \nu_1-\nu'_1,\nu_2-\nu'_2\ldots\nu_n-\nu'_n,
\]
that is different from zero has a \gesperrt{positive} value, that is, when either $\nu_1>\nu'_1$ or $\nu_1=\nu'_1,\nu_2>\nu'_2$ or $\nu_1=\nu'_1,\nu_2=\nu'_2,\nu_3>\nu'_3$, and so on.

Since we suppose that all terms in which all the exponents $\nu_1,\nu_2\ldots\nu_n$ agree have been combined into one term, this decides which of any two terms is higher; and if $A$ is higher than $A'$, and $A'$ higher than $A''$, then $A$ is also higher than $A''$.

If now, according to this ordering,
\[
 A=M\alpha_1^{\nu_1}\alpha_2^{\nu_2}\ldots\alpha_n^{\nu_n}
\]
is the highest term of our function $S$, it follows that
\[
 \nu_1\eqslantgtr\nu_2
\]
must hold; for if $\nu_1<\nu_2$, the term
\[
 M\alpha_1^{\nu_2}\alpha_2^{\nu_1}\ldots\alpha_n^{\nu_n},
\]
which must also occur in $S$ because of symmetry, would stand higher in the order than $A$, contrary to the supposition. Likewise, $\nu_2\eqslantgtr\nu_3$ must hold; for if $\nu_2<\nu_3$, then
\[
 M\alpha_1^{\nu_1}\alpha_2^{\nu_3}\alpha_3^{\nu_2}\ldots\alpha_n^{\nu_n}
\]
would stand higher than $A$, and so forth. We therefore conclude that the exponents in $A$
%CONT>>
